\chapter{LITERATURE SURVEY}

\section{Overview}
The development of autonomous robotic ground vehicles requires the convergence of computer vision, embedded real-time systems, machine learning, and control theory. Early automated guided vehicles relied on rigid physical modifications, such as magnetic tracks or optical tape detected via phototransistors, rendering them incapable of adapting to unstructured environments. With modern microprocessors, contemporary robotics shifted toward on-board vision-based perception. However, executing vision algorithms on mobile platforms introduces computational and electrical challenges. Research in this field focuses on optimizing neural networks for embedded runtimes, developing robust computer vision for lane extraction, designing multi-sensor collision avoidance, and building low-latency communication frameworks.

\section{Deep Learning and Quantization on Embedded Edge Systems}
Deep Convolutional Neural Networks (CNNs) have set state-of-the-art benchmarks in object classification, semantic segmentation, and bounding-box object detection. Single-stage object detectors, particularly the You Only Look Once (YOLO) family originally developed by Redmon et al., consolidated region proposal and bounding-box regression into a unified, end-to-end differentiable feedforward network. While earlier architectures like YOLOv3 and YOLOv4 required desktop-grade graphics processing units (GPUs) consuming hundreds of watts, recent lightweight architectures such as YOLOv8 Nano (YOLOv8n) were specifically designed for edge devices.

Despite these algorithmic advances, running full 32-bit floating-point (FP32) weights on ARM-based single-board computers like the Raspberry Pi results in unacceptable thermal throttling and frame processing times exceeding 300--400~ms. To address this limitation, post-training quantization techniques have been widely investigated. By mapping 32-bit floating-point weights and activation tensors to 8-bit integers (INT8) using symmetric uniform affine quantization, model storage is reduced by up to 75\%, while leveraging ARM NEON Single Instruction Multiple Data (SIMD) instruction sets. Modern inference engines such as Google LiteRT (formerly TensorFlow Lite) execute INT8 quantized models with negligible loss in Mean Average Precision (mAP) while delivering significant speedups. Recent studies confirm that deploying INT8-quantized YOLO models on ARM Cortex-A72 processors reduces single-frame inference latency down to approximately 100~ms, enabling practical on-board inference for micro-vehicles without cloud dependency.

\section{Classical Computer Vision for Lane Tracking and Corridor Clearance}
While deep neural networks excel at high-level semantic object recognition, classical computer vision remains the preferred paradigm for deterministic, low-latency road and lane tracking due to its low computational overhead. Traditional lane detection pipelines typically employ edge detection operators such as Canny filtering followed by Hough transforms to identify linear roadway boundaries. However, Hough transforms are computationally demanding and highly sensitive to lighting variations, shadow artifacts, and broken track markings.

Recent embedded robotics implementations emphasize centroid-based contour extraction over pre-filtered polygonal Regions of Interest (ROI). By applying Gaussian smoothing and Otsu or adaptive thresholding to isolate binary track markings, the spatial centroid $(\bar{x}, \bar{y})$ of the largest segmented contour can be calculated directly from image spatial moments:
\begin{equation}
\bar{x} = \frac{M_{10}}{M_{00}}, \quad \bar{y} = \frac{M_{01}}{M_{00}}
\end{equation}
The horizontal offset between the contour centroid $\bar{x}$ and the frame midpoint represents the lateral tracking error. This error serves as the proportional input to a discrete Proportional-Integral-Derivative (PID) steering controller. Furthermore, to handle obstacles unrepresented in standard training sets (such as plain walls or generic cardboard barriers), researchers have implemented corridor density scanning. By dividing the lower forward driving region into angular sectors and computing edge and pixel gradients, vehicles can measure clear space and safely identify evasion corridors in real time.

\section{Multi-Sensor Collision Avoidance and Distance Sensing}
A known vulnerability of purely monocular vision systems is the absence of direct depth information. Monocular depth estimation using deep neural networks is computationally prohibitive on low-power single-board computers, while looming-based Time-to-Collision (TTC) calculations based on bounding-box dilation rates $\frac{d(\text{area})}{dt}$ can be susceptible to optical noise and sudden camera pitch vibrations.

Consequently, multi-sensor fusion architectures combining monocular vision with direct active ranging sensors have emerged as the standard design pattern in robust mobile robotics. Ultrasonic sensors, such as the HC-SR04, measure Time-of-Flight (ToF) acoustic echo reflections at 40~kHz. While ultrasonic sensors suffer from limited beam divergence angles (typically $15^\circ$) and specular reflection at oblique incident angles, their direct physical distance measurements are unaffected by poor optical illumination, transparent glass, or lack of visual texture. Studies in autonomous micro-vehicles indicate that employing ultrasonic sensors as a deterministic, hard real-time safety boundary (e.g., triggering emergency braking at distances below 15~cm) provides an essential failsafe mechanism that prevents collisions when optical perception models suffer inference lag or false-negative detections.

\section{Real-Time Teleoperation Protocols and Telemetry Architectures}
In autonomous vehicle development, reliable human-in-the-loop teleoperation, remote diagnostic telemetry, and wireless fail-safes are critical for experimental safety. Traditional robotic architectures frequently employed heavy Robot Operating System (ROS) topic bridges, which introduce substantial operating system overhead and memory footprints that can strain micro-vehicles running on single-board computers.

Recent embedded designs favor lightweight web technologies. Micro-frameworks such as Flask, combined with Motion-JPEG (MJPEG) multipart streaming, enable non-blocking distribution of annotated camera video directly to standard HTML5 web browsers without client-side plugins. Furthermore, asynchronous RESTful endpoints permit bidirectional telemetry exchange (commanded throttle, steering angles, ultrasonic ranges, and FSM state) at sub-50~ms latencies. To safeguard against Wi-Fi disconnections, state-of-the-art systems implement dual wireless stacks, integrating Bluetooth RFCOMM SPP as an independent physical override channel alongside software watchdog auto-stop timers.
